\section{評価と考察}
\label{evaluation}
\ref{result}節で示した結果について評価と考察を行う．
\subsection{光通信の評価と考察}
光通信における動作結果，実験Bの動作結果と目標を表\ref{table:light}に示す．
通信精度，速度ともに目標に到達した．
この結果からシステムが有効かを評価する．
\par
実験では，LEDの点灯時間を少しずつ短くしながら行った．
20 msから短くしていき，5msより短くなると最後のシンボルを二重に受け取ることが増え，FSRが0.9を下回ってしまった．
実験Bでも同様の誤りが発生した．
5 msで点灯，消灯の検知を行うのが照度センサの限界であり，このシステムの最大限の速度で通信が行えた．
また，実験Bの結果，FSRははじめの実験よりも向上し，システムのより詳細な精度が得られた．
以上より，光通信は有効であると考える．

\par
このような結果が得られた原因について考察する．
本システムにおいて，光通信はLEDと照度センサに覆うようにカゴを設置して，暗い環境で行った．
これはカゴが無いと照明の影響を受けてしまい，照度センサが光の検出ができなかったためである．
また，LEDの明るさと照度センサの反応の強度のバランスを取ると，通信距離が5 cmとなった．
これらの光を検知するのに適切な環境設置をしたことが，精度の目標に到達した原因だと考える．
また，音通信ではGIを設けたが，光通信ではこのような信号の間隔は開けなかった．
これによって，必要最低限の信号の通信を行うことができ，速度が目標を超えたのだと考える．
LEDと照度センサの向きを変えて実験したり指向性が強いレーザー光を使用したりすることで通信距離の最大化が図れるだろう．


\begin{table}[tb]
    \centering
    \caption{光通信における動作結果と目標}
    \label{table:light}
    \begin{tabular}{|c|c|c|} \hline
            &通信精度  &通信速度  \\ \hline
        目標 & $FSR \geq 0.90$ &100 bps以上 \\ \hline
        動作結果 &$FSR = 0.90$ &143 bps  \\ \hline
        動作結果（実験B） &$FSR = 0.94$ &- \\ \hline
    \end{tabular}
\end{table}

\subsection{音通信の評価と考察}
音通信における動作結果と目標を表\ref{table:sound}に示す．
通信精度，速度ともに目標には至らなかった．
この結果からシステムが有効かを評価する．
% \par
% このような結果が得られた原因について考察する．
% 精度について，目標に到達はできなかったが，誤差は0.01であった．
% \ref{target}節で述べた通り，類似研究では通信音量より雑音が大きいときにBERが10 \%（0.10）だった．
% 今回の実験においても，他のグループの実験や会話があり，環境音はかなり影響があったと思われる．
% これらがスピーカからの音と干渉して正しい周波数を検出できなかったのだと考えられる．
% 実験時に環境音と通信音量をそれぞれ測っておくことでこの問題の解決につながるだろう．
% また，特定の音だけ抜ける事があり，周波数によって検出しやすい・しにくい距離があったように考えられる．
% 2000 $\thicksim$ 3600 Hzの1600 Hzの幅で同じように検出できる距離が見つからなかったことも原因だろう．
このBERで5回連続で同じ文字列を送信することを考えると，5回すべて誤る確率は$0.11^5 \simeq 0.00002$であり，少なくとも1回は完全な送信ができる確率は0.99998となる．
目標のBERでも同様にすると，少なくとも1回は完全な送信ができる確率は0.99999であり，目標の99.998 \%の数値となった．
また，1回あたりの通信時間は約5秒であり，本システムは25秒あれば99.998 \%の正確な通信ができる事となる．
精度については有効なシステムであると考える．
しかし，速度については目標の67 \%の数値であることから，全体で目標の約67 \%の性能であるといえる．
音通信全体として考えると精度は申し分ないが速度が足りず，結論として有効なシステムではないと考える．
\par
このような結果が得られた原因について考察する．
はじめに精度について考察する．
今回の実験では他のグループの実験や会話があり，環境音はかなり影響があったと思われる．
これらがスピーカからの音と干渉して正しい周波数を検出できなかったのだと考えられる．
また，特定の音だけ抜ける事があり，周波数によって検出しやすい・しにくい距離があったように考えられる．
2000 $\thicksim$ 3600 Hzの1600 Hzの幅で同じように検出できる距離が見つからなかったことも原因に挙げられる．
実験時に環境音と通信音量をそれぞれ測っておき，その割合ごとに実験することでこの問題の解決につながるだろう．
また，誤り検出や再送処理を導入することでより精度の高いシステムになるといえる．
\par
次に速度について考察する．
同期信号を1回ずつ出すところを急遽2回ずつ出すように変更したことや，終了の同期の安定のためにGIを1つ増やしたことなどが目標に至らなかった原因として考えられる．
これらをカバーするために音長を少しずつ短くして実験できれば，速度の改善が見られるだろう．
また，環境音のノイズを軽減する処理もシステムの有効化に必要である．

% \par
% 別の課題として，Arduinoのtone関数を用いてスピーカから出る音があまり聞き心地が良くなかったことも挙げられる．
% 通信を可聴化するということにおいて，聞き心地が良くないと通信の快適さが精神的な面で損なわれてしまう．
% 仕組みは複雑になるが，他の音と組み合わせて和音にするなどの工夫が必要である．

\begin{table}[tb]
    \centering
    \caption{音通信における動作結果と目標}
    \label{table:sound}
    \begin{tabular}{|c|c|c|} \hline
            &通信精度  &通信速度  \\ \hline
        目標 & $BER < 0.10$ &10 bps以上 \\ \hline
        動作結果 &$BER = 0.11$ &6.7 bps\\ \hline
    \end{tabular}
\end{table}

\subsection{振動通信の評価と考察}
振動通信における動作結果と目標を表\ref{table:vibration}に示す．
通信精度，速度ともに目標には至らなかった．
\par
このような結果が得られた原因について考察する．
精度について，結果は目標の20倍以上と大きく下回ってしまった．
振動モータを高速でオンオフさせると，1つ前の振動が残ってしまい，これが誤検知につながっている．
これに加えて，加速度センサ自身が振動モータの振動の影響を受けてしまっているのでは無いかと考えた．
次に速度について，結果は目標の30 \%以下となった．
振動でデータを送受信するとなると，データの表現方法に乏しく，4進数に変化にして送信するといったようなデータの短縮や効率化が他のシステムに比べ困難である．
振動モータの出力ではセンサに加える加速度が安定せず，これにより誤検知を軽減するために1回あたりの振動時間を長くしたことが原因だと考えた．
また，課題として閾値を一度設定しても，土台を乗せる環境ごとに精度が変化し再設定する必要があり，どんな環境でも利用できるわけではないことが挙げられた．


\begin{table}[tb]
    \centering
    \caption{振動通信における動作結果と目標}
    \label{table:vibration}
    \begin{tabular}{|c|c|c|} \hline
            &通信精度  &通信速度  \\ \hline
        目標 & $BER < 0.10$ &10 bps以上 \\ \hline
        動作結果 &$BER = 0.20$ &2.81 bps\\ \hline
    \end{tabular}
\end{table}

\subsection{システム全体の評価と考察}
システム全体の目標として，精度については各部門におけるbitの一致率を90 \%として，これらの期待値$0.90^{3} = 0.729$を設定する．
速度について，各部門における目標bpsが光で100，音，振動で10であり，1 bitあたりの送信時間はこれらの逆数である．
逆数の和から再度逆数を取ることで，全体の速度の目標を4.76 bps以上を設定する．
\par
システム全体の精度は全10回の試行で誤りビット数が88 bit，総送信ビット数が350であるため，$BER \simeq 0.2514$となる．
よって，ビット一致率は$(1-BER)\times100 = 74.86 \%$である．
ビット一致率とは別に文字としても一致率も求める．
不一致の文字数が34，総文字数が50であるため，文字不一致率が68 \%，つまり文字一致率は32 \%である．
速度については総送信ビット数が350 bit，総通信時間の平均が207.58 sであるため，1.69 bpsとなる．

\par
このような結果が得られた原因について考察する．
精度については目標を上回ることができた．
しかし，ビットが70 \%一致していてもデコードされた文字をみて"chiba"が送信されたとは思えないだろう．
ここで文字一致率をみると32 \%であり，例えば"chiba"を送受信してデコードされた文字列が"dhkba"といったようになる．
これは最後に2進数に変換していることが原因だと考える．
1文字あたり，4進数4桁のとき，2進数では8桁が必要になり，1桁あたりの重みは2進数の方が小さいため，1桁誤ったときのビット（シンボル）一致率は2進数のほうが高くなる．
しかし，どちらも1桁誤ると異なる文字にデコードされるもしくは文字化けすることは共通であり，桁数が多い分ビット一致率が高くても文字の一致率は2進数の方が低くなると考える．
次に，速度について，目標の約30 \%となった．
光に比べて音と振動の速度がかなり遅く，これが速度の低下の原因である．
本システムでは精度を向上させるために速度を抑えたが，これで文字一致率が32 \%であることを考えると，通信速度に重点をおいてシステムを構築するほうが有効なシステムになると考えた．
これらのことから，本システムは文字を伝送するシステムとしての有効性は低いと言える．